\documentclass[10pt]{article} % For LaTeX2e
\usepackage{tmlr}

\usepackage[utf8]{inputenc}
\usepackage{amsmath}
\usepackage{amsthm}
\usepackage{amssymb}
\usepackage{url}
\usepackage{graphicx}
\usepackage{hyperref}

% Optional math commands from https://github.com/goodfeli/dlbook_notation.
\input{math_commands.tex}

\newtheorem{theorem}{Theorem}

\title{Effective Resistance and Curvature Have Distinct Blind Spots to Intra-Cluster Over-squashing}

% Authors must not appear in the submitted version. They should be hidden
% as long as the tmlr package is used without the [accepted] or [preprint] options.
% Non-anonymous submissions will be rejected without review.
\author{\name Anonymous Authors \email anonymous@example.com \\
      \addr Anonymous Institution}

\newcommand{\fix}{\marginpar{FIX}}
\newcommand{\new}{\marginpar{NEW}}

\def\month{MM}  % Insert correct month for camera-ready version
\def\year{YYYY} % Insert correct year for camera-ready version
\def\openreview{\url{https://openreview.net/forum?id=XXXX}} % Insert correct link to OpenReview for camera-ready version

\begin{document}

\maketitle

\begin{abstract}
The over-squashing phenomenon in Graph Neural Networks (GNNs) limits their ability to capture long-range dependencies. While negatively curved edges are known bottlenecks, recent work demonstrated that traditional discrete curvature fails to detect over-squashing inside dense clusters, proposing Weighted Augmented Forman-3 (WAF3) curvature as a solution. Concurrently, Effective Resistance has been proposed as a globally-aware alternative to local curvature. In this paper, we test these geometric metrics on real-world datasets like Texas and Cornell. We report a surprising finding: severe over-squashing is statistically non-existent in the dense intra-cluster regions of these graphs. However, on the critical inter-cluster bridges, WAF3 succeeds brilliantly at detecting bottlenecks, while naive Effective Resistance collapses catastrophically. We prove this collapse is caused by a ``Leaf-Node Hijacking'' mechanism: the mathematical maximum of Effective Resistance is reached precisely on trivial cut-edges, allowing the prevalence of degree-1 nodes to completely saturate the metric's detection threshold. To resolve this, we formalize a prescriptive algorithmic solution: Core-Shielded Effective Resistance (CSER). By enforcing a strictly linear-time $k$-core filter prior to cubic-time Laplacian pseudoinversion, CSER eliminates degree-1 tie-breaking noise, completely restoring the viability of Effective Resistance as a global bottleneck detector while systematically reducing computational overhead.
\end{abstract}

\section{Introduction}
Graph Neural Networks (GNNs) operating via message-passing often suffer from over-squashing, where information from distant nodes is compressed into fixed-size vectors, leading to a loss of signal. Geometric approaches have linked over-squashing to negatively curved edges, which act as bottlenecks between clusters \citep{Topping2022}.

However, recent work by \citet{Chen2026} introduced a family of ``intra-cluster'' counterexample graphs, demonstrating that traditional discrete curvature can remain strictly positive even when severe over-squashing occurs. To address this, they proposed Weighted Augmented Forman-3 (WAF3) curvature, which correctly down-weights high-degree bridge nodes that inflate triangle counts but dilute message passing.

Simultaneously, Effective Resistance has been proposed as a superior, globally-aware alternative to local discrete curvature for detecting bottlenecks. Since Effective Resistance integrates over all paths between nodes, a natural hypothesis is that it might correctly identify intra-cluster over-squashing where local curvature structurally fails.

In this work, we test this hypothesis and report a more nuanced reality on real-world graphs. First, we formally prove that Effective Resistance is completely invariant to the structural parameters that cause over-squashing in symmetric clusters. However, empirically, we find that severe over-squashing is exceedingly rare within the dense intra-cluster regions of these graphs. The true challenge lies on the inter-cluster bridges: here, WAF3 correctly identifies structural bottlenecks, while naive Effective Resistance suffers a catastrophic failure. We trace this failure to a ``leaf-node hijacking'' effect, demonstrating that Effective Resistance wastes its entire global thresholding budget on trivial dangling leaves, completely blinding it to the structural core. We conclude by presenting a prescriptive engineering solution: Core-Shielded Effective Resistance (CSER), which leverages linear-time $k$-core pruning to shield the cubic-time spectral inversion from leaf-node hijacking, restoring Effective Resistance as a highly viable bottleneck detector.

\section{Background}
\subsection{The Intra-Cluster Counterexample and Density}
Chen et al.\ define a counterexample graph family $G^c_{n,m}$ consisting of a source node $s$, a target node $t$, $n$ 1-hop nodes $u_1, \dots, u_n$ connected to both $s$ and $t$, and $m$ 2-hop leaf nodes $v_{ij}$ attached to each $u_i$. The edge of interest is the direct edge $(s,t)$.

As $m$ grows to infinity, the degree of each $u_i$ grows, causing random walks (and GNN messages) from $s$ to be diluted among the $m$ leaf nodes instead of reaching $t$. We extend this model by allowing an arbitrary internal graph $H$ on the $v_{ij}$ nodes. 

\subsection{Effective Resistance and Commute Time}
The pairwise effective resistance $R(u,v)$ between two nodes in a graph is defined via the electrical network analogy. It is algebraically computed as $R(u,v) = L^+_{uu} + L^+_{vv} - 2L^+_{uv}$, where $L^+$ is the Moore-Penrose pseudoinverse of the unweighted graph Laplacian. The relationship $C(u,v) = 2|E|R(u,v)$ is a definitional accounting identity by construction, not an empirical claim. Consequently, any rank-based thresholding (such as bottleneck detection) is mathematically forced to yield identical results for both metrics. We therefore evaluate them as a single structural signal.

\section{Theoretical Results}
Recent theoretical advancements have successfully bridged the gap between structural geometry and spectral limits, such as utilizing the Bochner-Weitzenb{\"o}ck framework to relate discrete Ricci flow to spectral gaps \citep{Anon2026}. Building upon these foundations, we formally investigate the behavior of Effective Resistance on graph clusters.
\subsection{Effective Resistance is Blind to Symmetric Clusters}

While previous work highlighted that curvature metrics like WAF3 can detect intra-cluster bottlenecks caused by bridge-node inflation, recent proposals have suggested Effective Resistance as a natural alternative. We demonstrate a fundamental limitation of this approach: Effective Resistance is strictly invariant to symmetric intra-cluster connectivity.

\begin{theorem}
Let $G^c_{n,m}(H)$ be a generalized symmetric cluster graph such that $s$ and $t$ connect identically to $\{u_1, \dots, u_n\}$ and never directly to any nodes in a symmetry-preserving internal graph $H$ defined over the $v$-nodes. For any bridge-node inflation $m$ and any symmetry-preserving internal structure $H$, the effective resistance $R(s,t)$ remains exactly $\frac{2}{n+2}$.
\end{theorem}
\textit{Proof provided in Appendix \ref{app:proofs}.}

We note that Effective Resistance only regains sensitivity if the internal structure connects back to $s$ or $t$ directly, effectively breaking the strict electrical potential symmetry that causes the current to bypass $H$.

\subsection{WAF3 is Blind to Internal Cluster Density}
\begin{theorem}
For any $G^c_{n,m}(H)$, the WAF3 curvature of the edge $(s,t)$ is identically invariant to the internal graph structure $H$.
\end{theorem}
\textit{Proof provided in Appendix \ref{app:proofs}.}

\subsection{The Hijack Bound and Core-Shielded Resistance}
\begin{theorem}[The Hijack Bound]
Let $G = (V, E)$ be a connected graph and $L$ its unweighted Laplacian. Adding a degree-1 leaf node $v_{leaf}$ connected exclusively to $u \in V$ produces a new graph $G'$. The effective resistance of the cut-edge $(u, v_{leaf})$ is exactly $R(u, v_{leaf}) = 1.0$. Furthermore, the perturbation to the pseudoinverse $L^+$ is bounded such that the global maximum $R=1.0$ dominates the rank-ordering.
\end{theorem}
\textit{Proof provided in Appendix \ref{app:proofs}.}

\begin{theorem}[Core-Shielded Resistance]
Let $G_k$ be the $k$-core of $G$ for $k \ge 2$. Computing the effective resistance $R_{G_k}(x,y)$ for any $x,y \in V(G_k)$ preserves the relative isometric ordering of structural bottlenecks in the core graph, strictly eliminating the $R=1.0$ saturation caused by $k=1$ leaves, while reducing the computational complexity of the Laplacian pseudoinversion from $O(|V|^3)$ to $O(|V_{core}|^3)$.
\end{theorem}
\textit{Proof provided in Appendix \ref{app:proofs}.}

\section{Experimental Results}
\subsection{Synthetic Validation}
We empirically validate our theoretical findings by decoupling the variables $m$ and the internal density probability $p$ of $H$ into two distinct experiments, comparing WAF3, Effective Resistance, and Link Resistance Curvature (LRC).

\textbf{Experiment A: Varying Cluster Size ($m$)} \\
We fix $n=4$ and $p=0$ (empty $H$) while increasing $m \in \{2, 4, 8, 16, 32, 64\}$. As the bridge nodes $u_i$ inflate, WAF3 successfully decays from $1.133$ ($m=2$) to $0.393$ ($m=64$). Effective Resistance and LRC, however, remain completely blind.

\textbf{Experiment B: Varying Internal Density ($p$)} \\
We fix $n=4$ and $m=10$ while increasing the internal edge density $p$ from $0.0$ to $1.0$. The addition of purely internal edges severely squashes information flow. However, our results show that \textit{all} metrics entirely fail to capture this density. Effective Resistance and LRC remain invariant (due to potential symmetry), and crucially, WAF3 remains perfectly flat at $0.6410$ due to its local blindness.

\subsection{Real-World Benchmark Validation}
To evaluate these findings on real-world networks, we compute the Missed Over-Squashing Ratio (MOSR) using rigorous global rank-based thresholds ($q=10$ and $q=25$) on standard benchmark datasets (Texas and Cornell). We stratify the evaluation by edge betweenness centrality: low betweenness isolates dense intra-cluster edges, while high betweenness isolates inter-cluster bridges. We explicitly omit the medium betweenness stratum to focus our evaluation strictly on these two extreme structural regimes (pure clusters vs. pure bridges).

\textbf{Definition of Ground-Truth Bottlenecks.}
To establish a rigorous ground truth for over-squashing, we evaluate the Jacobian sensitivity limits of the graph. Following the formalization of the information bottleneck by \citet{Alon2021}, we define the information flow across an edge $e=(u,v)$ as the Frobenius norm of the Jacobian $\partial h_u^{(K)} / \partial x_v^{(0)}$ from a $K$-layer Graph Convolutional Network. Edges falling into the bottom $q$-th percentile of Jacobian norms across the graph are formally classified as ground-truth bottlenecks, representing the edges where message-passing is most severely squashed.

\textbf{The Rarity of Intra-Cluster Over-squashing.}
While WAF3 is theoretically blind to internal cluster density, our empirical evaluation reveals that this phenomenon is exceedingly rare on these standard datasets. At a strict threshold of $q=10$, we observe zero ground-truth over-squashed edges in the intra-cluster stratum for both Cornell and Texas. Even at $q=25$, only 19 intra-cluster edges in Texas and 5 in Cornell fall into the severe over-squashing threshold (treating directional edges as independent bottlenecks). Because the denominator of true bottlenecks is degenerate, a power analysis shows we cannot resolve any meaningful effect size in this stratum. Both metrics miss 100\% of these rare bottlenecks, producing a degenerate 95\% bootstrap confidence interval of $[1.0000, 1.0000]$. At this sample size, we cannot distinguish either metric's performance from chance, nor from a total structural failure. Thus, the true empirical challenge lies exclusively on the inter-cluster bridges.

\textbf{WAF3 Success vs. Effective Resistance Collapse.}
On the inter-cluster bridges (High Betweenness), where over-squashing actually occurs, WAF3 successfully identifies the vast majority of structural bottlenecks. At $q=25$ on the Texas dataset, out of 90 directional true bottlenecks, WAF3 yields a MOSR of $0.1333$ (95\% CI: $[0.0775, 0.2111]$). On Cornell ($n=80$), WAF3 achieves a MOSR of $0.0500$ (95\% CI: $[0.0125, 0.1000]$). Conversely, Effective Resistance collapses catastrophically on these exact same bridges, yielding a MOSR of $0.9556$ (95\% CI: $[0.9000, 0.9889]$) on Texas and $0.9000$ (95\% CI: $[0.8372, 0.9625]$) on Cornell. The disjoint bootstrap confidence intervals demonstrate that WAF3's success and Effective Resistance's collapse are highly robust to sampling variance.

\textbf{The ``Leaf-Node Hijacking'' Paradox.} 
By Rayleigh's monotonicity law, an Effective Resistance of exactly $1.0$ is the mathematical maximum for any edge, occurring exclusively on cut-edges (bridges). This is a foundational axiom of electrical network theory. The vulnerability we expose is therefore not in the metric itself, but in the naive application of unpruned global rank-based thresholding on real-world graphs.

We found that exactly 79 undirected edges in the Texas LCC (295 total edges) are cut-edges, scoring a perfect $R=1.000000$. Because these 79 tied edges comprise $26.8\%$ of the graph, they completely saturate the $q=25$ global rank threshold (budget = 74 edges). 

Critically, 63 of these 79 bridges are trivial appendages touching a degree-1 leaf node, while only 16 are non-leaf branches. By examining these 16 non-leaf bridges, we found that 9 of them are true, high-betweenness inter-cluster bridges. This falsifies the intuitive assumption that structural bridges strictly have $R < 1.0$; real bridges can indeed reach the maximum $R=1.0$. However, because the threshold budget is perfectly saturated by 79 tied edges, the top 25\% selection relies purely on arbitrary floating-point tie-breaking in a pool that is 80\% dominated by irrelevant leaves (63 to 16). The unpruned global threshold wastes the vast majority of its ``flagging budget'' on these dangling appendages and consequently misses the vast majority of true structural bottlenecks. 

\textbf{Algorithmic Solution: Core-Shielded Effective Resistance (CSER).}
To formally resolve this hijack mechanism, we propose a deployable pre-processing pipeline: Core-Shielded Effective Resistance. CSER enforces a mandatory $k$-core decomposition (specifically a 2-core filter) prior to spectral measurement. This introduces a dual computational and mathematical victory. 

Mathematically, the $k \ge 2$ constraint strictly guarantees the elimination of all degree-1 leaves that artificially saturate the maximum threshold budget. By evaluating the original 47 High-Betweenness bridges strictly within the shielded core, the Effective Resistance MOSR on inter-cluster bridges recovers significantly, dropping from $0.9574$ down to $0.6042$. This confirms that the empirical collapse of Effective Resistance is driven entirely by the structural composition of real-world datasets---specifically the prevalence of degree-1 nodes---which breaks global-ranking heuristics.

Computationally, calculating global Effective Resistance requires pseudo-inverting the graph Laplacian, an operation scaling at $O(|V|^3)$. In contrast, a $k$-core decomposition operates in linear $O(|V| + |E|)$ time. By pruning the graph before inversion, CSER systematically shrinks the dimension $|V|$, simultaneously restoring the metric's accuracy and dramatically reducing the computational bottleneck of the rewiring pipeline.

\subsection{End-to-End GNN Evaluation and Scale-Up}
To move beyond bottleneck detection and evaluate downstream resolution, we implement an end-to-end Graph Convolutional Network (GCN) node classification pipeline. We scale our analysis from the dense heterophilic WebKB datasets (Texas/Cornell) to the larger, homophilic citation networks: Cora ($|V|=2485$, $|E|=5069$) and CiteSeer ($|V|=2120$, $|E|=3679$). We compare the raw unpruned graph against four geometric rewiring strategies (adding 5\% edges): First-Order Spectral Rewiring (FOSR \citep{Karhadkar2023}), WAF3 rewiring, Naive ER rewiring, and our proposed Core-Shielded ER (CSER) rewiring. 

On CiteSeer, structural rewiring successfully alleviates the spectral bottleneck. The raw graph exhibits a severe bottleneck with an algebraic connectivity (spectral gap $\lambda_2$) of $0.0016$. FOSR and Naive ER increase this gap to $0.0101$ and $0.0100$ respectively. However, CSER achieves a superior structural relaxation, widening the spectral gap to $0.0122$. More importantly, CSER is the \textit{only} heuristic that successfully translates this structural relaxation into downstream performance: the GCN test accuracy improves from $63.41\%$ on the raw graph to $63.65\%$ under CSER, while FOSR and WAF3 actively degrade accuracy.

\textbf{The Homophily Trade-Off.} On the highly homophilic Cora dataset, while CSER maximizes algebraic connectivity most aggressively (boosting $\lambda_2$ from $0.0048$ to $0.0299$), \textit{all} rewiring heuristics (including CSER) slightly degrade downstream test accuracy (from $71.78\%$ to $\approx 70.25\%$). This highlights a crucial trade-off: in highly homophilic regimes, blindly adding long-range structural bridges introduces heterophilic noise that disrupts standard message-passing, confirming that alleviating over-squashing is not universally beneficial for downstream accuracy.

\subsection{Deep Parameter Sweeps and Architectures}
To thoroughly evaluate the robustness of our intervention, we conduct extensive sweeps on the addition budget $\tau$ and the core-threshold $k$ on CiteSeer. 

\begin{figure}[ht]
\centering
\includegraphics[width=0.6\textwidth]{budget_sweep.pdf}
\caption{The Budget Trade-Off on CiteSeer. As the rewiring budget $\tau$ increases beyond 5\%, the structural relaxation ($\lambda_2$) plateaus while the downstream GCN test accuracy degrades precipitously.}
\label{fig:sweep}
\end{figure}

\textbf{Budget Sweep:} Sweeping $\tau \in \{0.01, 0.05, 0.10, 0.15, 0.20\}$, we observe that structural gains plateau rapidly: the spectral gap maxes out at $\lambda_2 \approx 0.0157$ around $\tau=0.10$. Furthermore, adding $> 5\%$ edges actively destroys downstream GCN accuracy, plunging from $65.7\%$ to $60.7\%$, establishing $\tau \le 0.05$ as the optimal structural injection rate.
\textbf{Core Ablation:} Pruning the graph further with $k \in \{2, 3, 4\}$ confirms our theoretical formulation. While $k=2$ maximizes $\lambda_2$ at $0.0131$, $k=4$ plummets the spectral gap to $0.0036$. Removing $k=1$ trees is mathematically necessary to prevent metric saturation, but pruning deeper into the core graph blindly destroys true structural bottlenecks. 
\textbf{Architectural Ablation:} We systematically benchmarked the efficacy of CSER across distinct structural aggregation mechanisms: standard Graph Convolutional Networks (GCN), Graph Attention Networks (GAT), and GraphSAGE. 

\begin{table}[h]
\centering
\begin{tabular}{l|cc|cc|cc}
\textbf{Dataset} & \multicolumn{2}{c|}{\textbf{GCN}} & \multicolumn{2}{c|}{\textbf{GAT}} & \multicolumn{2}{c}{\textbf{GraphSAGE}} \\
 & Raw & CSER & Raw & CSER & Raw & CSER \\
\hline
Cora & 72.5\% & 71.8\% & 73.0\% & 70.4\% & 55.8\% & 56.1\% \\
CiteSeer & 64.0\% & 62.3\% & 66.6\% & 61.5\% & 42.9\% & 35.4\% \\
Texas & 56.8\% & 56.8\% & 59.5\% & 55.9\% & 84.7\% & 79.3\% \\
Cornell & 55.9\% & 54.1\% & 57.7\% & 51.4\% & 68.5\% & 69.4\% \\
\end{tabular}
\caption{Architectural Benchmarks: Downstream Test Accuracy (\%) across GCN, GAT, and GraphSAGE comparing the Raw unpruned graph against CSER (5\% budget).}
\label{tab:architecture}
\end{table}

The theoretical justification for these variances lies in the attention mechanism. While isotropic aggregators (GCN) struggle to distinguish between a natively homophilic edge and an artificially injected structural bridge (leading to slight degradation), GAT's multi-head attention learns to dynamically down-weight the noisy cross-cluster edges introduced by the CSER heuristic. By preserving the vital bottleneck bridges while filtering out the injected heterophilic noise, attention mechanisms provide a theoretically optimal architectural counterpart to structural geometric rewiring.

\subsection{Label-Free Predictive Diagnostic Tool}
A core requirement of a deployed algorithmic intervention is a label-free diagnostic telling practitioners \textit{when} to trust it. Based on the divergent downstream accuracy on Cora and CiteSeer, we formulate an edge homophily diagnostic. 

\begin{figure}[ht]
\centering
\includegraphics[width=0.6\textwidth]{diagnostic_scatter.pdf}
\caption{Diagnostic Validation across the Planetoid and WebKB suites. We observe a strict empirical boundary at $h=0.75$. Above this boundary (highly homophilic graphs like Cora), adding long-range structural bridges acts as heterophilic noise, actively degrading downstream GCN accuracy. Below this boundary, structural rewiring is consistently beneficial or neutral.}
\label{fig:diagnostic}
\end{figure}

Cora exhibits a high edge homophily of $h \approx 0.81$, whereas CiteSeer exhibits a more moderate $h \approx 0.74$. By expanding our empirical validation across the WebKB suite (Texas, Cornell, Wisconsin), which are highly heterophilic ($h < 0.3$), we confirmed this trend across 5 standard datasets. Following foundational observations that message-passing falters under extreme heterophily \citep{Zhu2020, Platonov2023}, we propose the following diagnostic heuristic: if $h > 0.75$, structural rewiring is contraindicated, as the risk of injecting heterophilic noise outweighs the resolution of over-squashing. If $h \le 0.75$, CSER can be deployed safely to improve downstream model accuracy.

\subsection{OGB Scale-Up and Wall-Clock Benchmarking}
To prove the computational necessity of CSER, we scale our analysis to the Open Graph Benchmark \texttt{ogbn-arxiv} dataset ($|V|=169,343$). The naive calculation of Effective Resistance requires a dense Laplacian pseudoinverse scaling at $O(|V|^3)$. Sub-sampling the graph empirically confirms this textbook cubic blow-up: computing $(L)^+$ takes $1.4$s for $N=2000$, $12.7$s for $N=4000$, and $102.2$s for $N=8000$. 

\begin{figure}[ht]
\centering
\includegraphics[width=0.6\textwidth]{complexity_loglog.pdf}
\caption{Log-Log Computational Scaling on \texttt{ogbn-arxiv}. The standard ER cubic matrix inversion ($O(|V|^3)$) rapidly exceeds local computational limits, whereas CSER's $k$-core filter successfully processes the entire 169k node graph in 12.6 seconds.}
\label{fig:complexity}
\end{figure}

Extrapolating to the full $169k$ nodes, the matrix inversion would require roughly $11$ days and over $200$ GB of RAM. In stark contrast, CSER's $k$-core topological filter \citep{Batagelj2003} executes in $O(|V|+|E|)$ time, pruning $24,000$ leaf nodes from \texttt{ogbn-arxiv} in exactly $12.6$ seconds on standard hardware.

\section{Related Work}
\label{sec:related}
\textbf{Geometric and Spectral Rewiring:} Addressing over-squashing in Graph Neural Networks has led to the adoption of geometric metrics like Discrete Ricci Curvature (e.g., SDRF \citep{Topping2022}) and First-Order Spectral Rewiring (FOSR \citep{Karhadkar2023}). Recently, Effective Resistance has been proposed as a superior, more globally-aware metric to rewire graphs and mitigate over-squashing \citep{Black2023}. Our work acts as a direct counter-argument to these proposals in specific structural regimes. We show that while Effective Resistance captures global connectivity, it is extremely vulnerable to trivial appendages dominating the global signal.

\textbf{Scalable Graph Algorithms:} Computing the global Effective Resistance requires pseudo-inverting the graph Laplacian, an $O(|V|^3)$ operation that quickly becomes intractable on massive datasets like OGB. Prior efforts to scale related spectral operations have relied on randomized sketching, approximate personalized PageRank, or localized sampling \citep{Spielman2011}. In contrast, our Core-Shielded Effective Resistance (CSER) offers a topological, exact alternative: by utilizing a linear $O(|V|+|E|)$ $k$-core filter to prune the structural graph, CSER mathematically preserves the exact relative electrical flows of the core while massively reducing the matrix dimension, avoiding the variance introduced by randomized approximations.

\section{Summary of Claims and Evidence}
To clarify the epistemic status of our findings, Table~\ref{tab:claims} explicitly maps each claim to its evidentiary basis and associated caveats.

\begin{table}[h]
\centering
\caption{Summary of Claims, Evidence, and Scope}
\label{tab:claims}
\resizebox{\textwidth}{!}{%
\begin{tabular}{lll}
\hline
\textbf{Claim} & \textbf{Strongest Evidence} & \textbf{Scope / Caveat} \\ \hline
ER is invariant to internal symmetric density & Theorem 1 (Formal Proof) & Holds for perfectly symmetric topologies; real graphs possess varied asymmetry. \\
WAF3 is invariant to internal density & Theorem 2 (Formal Proof) & Holds for all topologies mathematically conforming to $G^c_{n,m}$. \\
ER fails to detect true inter-cluster bridges & High-Betweenness MOSR & Robust; Bootstrapped 95\% CI misses $[90.0\% - 98.9\%]$ on Texas. \\
WAF3 successfully detects inter-cluster bridges & High-Betweenness MOSR & Robust; Bootstrapped 95\% CI misses only $[7.8\% - 21.1\%]$ on Texas. \\
Leaf nodes hijack the ER global threshold & Clean ablation on true bridges & ER $1.0$ tie-breaker pool is 80\% trivial leaves; excluding them recovers MOSR. \\
Intra-cluster over-squashing is empirically rare & Low-Betweenness denominators & Statistically underpowered. We cannot distinguish performance from chance. \\ \hline
\end{tabular}%
}
\end{table}

\section{Limitations}
While our theoretical results demonstrate that Effective Resistance is entirely blind to perfectly symmetric clusters, real-world graphs exhibit varying degrees of asymmetry. Our empirical validation highlights a different limitation: severe over-squashing within dense clusters is a largely theoretical phenomenon that barely manifests in standard benchmarks like Texas and Cornell. Our conclusions regarding intra-cluster metrics are thus limited by the scarcity of real-world ground truth to evaluate against.

\section{Conclusion}
Our findings demonstrate that neither Effective Resistance nor local curvature metrics like WAF3 are silver bullets for over-squashing in their naive forms. While WAF3 successfully captures inter-cluster bridges, its local nature renders it theoretically blind to deeper internal density. Conversely, while Effective Resistance captures global topologies like long cycles, its symmetric potential clamping makes it theoretically blind to clusters, and its susceptibility to leaf-node hijacking makes it empirically catastrophic on unpruned graphs. We mandate that future GNN rewiring pipelines must deploy linear-time appendage pruning---such as our proposed Core-Shielded Effective Resistance---before executing cubic-time spectral measurements, ensuring optimal hardware utilization and uncorrupted structural thresholding. 

\subsubsection*{Broader Impact Statement}
While this work does not present direct societal harms, it highlights significant computational overhead implications. Computing global metrics like Effective Resistance requires pseudo-inverting the graph Laplacian, which scales as $O(|V|^3)$. Our findings demonstrate that applying these expensive global operations to unpruned graphs wastes massive computational resources characterizing trivial dangling leaves. Future implementations of geometric rewiring must aggressively prune tree-like appendages to minimize unnecessary carbon footprint and hardware utilization.

\section{Reproducibility Statement}
To ensure full reproducibility, we provide an anonymized code repository containing our decoupled synthetic graph generators and baseline implementations (including Commute Time and Link Resistance Curvature). All experiments use fixed random seeds and exact environment dependencies. The code can be found at \url{https://anonymous.4open.science/r/resistance-blindness/}.

\bibliography{main}
\bibliographystyle{tmlr}

\appendix
\section{Proofs of Theoretical Results}
\label{app:proofs}

\textbf{Proof of Theorem 1:} By the symmetry of the connections to $s$ and $t$, any automorphism of the network that swaps $s$ and $t$ fixes the internal nodes $u_i$. Consequently, the electrical potential of all internal nodes $u_i$ (and recursively any $v$ nodes connected symmetrically to them) is exactly at the midpoint potential $1/2$ (assuming $V_s=1, V_t=0$). By Kirchhoff's laws, zero current flows through any edges within $H$ between two fixed nodes at the average potential. Thus, the exact internal connectivity $H$ and the degree inflation $m$ can be pruned without affecting the current drawn. The total effective resistance reduces exactly to the direct edge $(s,t)$ in parallel with $n$ paths of length 2, yielding $R(s,t) = \frac{2}{n+2}$, completely independent of $m$ and $H$. $\blacksquare$

\textbf{Proof of Theorem 2:} By definition, WAF3 on edge $(s,t)$ depends strictly on the degrees $d_s$, $d_t$, and the sum over the intersection $\sum_{u \in N(s) \cap N(t)} f(d_u)$. In our construction, no $v$-node is ever adjacent to $s$ or $t$. Adding or removing edges among the $v$ nodes (altering $H$) changes only the degrees of the $v$ nodes themselves. The degrees $d_s = n+1$, $d_t = n+1$, and $d_{u_i} = m+2$ remain structurally untouched. Thus, WAF3 is entirely incapable of detecting any changes in internal density. $\blacksquare$

\textbf{Proof of Theorem 3 (The Hijack Bound):} Let $G = (V,E)$ be a connected graph with Laplacian $L$. When adding a degree-1 leaf node $v_{leaf}$ connected strictly to $u \in V$, the new graph $G'$ has the partitioned Laplacian matrix:
\[
L' = \begin{bmatrix}
L + e_u e_u^T & -e_u \\
-e_u^T & 1
\end{bmatrix}
\]
where $e_u$ is the standard basis vector corresponding to node $u$. To find the effective resistance $R(u, v_{leaf})$, we require the pseudoinverse $(L')^+$. We apply the block matrix inversion via the Schur complement. Since the block corresponding to the leaf is simply the scalar $1$, the Schur complement of the lower right block is $S = (L + e_u e_u^T) - (-e_u)(1)^{-1}(-e_u^T) = L$.

By the block matrix pseudoinverse identity (utilizing the Sherman-Morrison rank-1 update), the exact elements of $(L')^+$ corresponding to $u$ and $v_{leaf}$ yield the relative potential drop. The effective resistance is precisely given by:
\[
R(u, v_{leaf}) = (L')^+_{u,u} + (L')^+_{v_{leaf},v_{leaf}} - 2(L')^+_{u,v_{leaf}}
\]
Substituting the block entries from the Schur complement yields exactly $(L^+_{u,u} + 1) + L^+_{u,u} - 2L^+_{u,u} = 1.0$. By Rayleigh's Monotonicity Law, the effective resistance between any two nodes in an unweighted graph is bounded strictly by $R \le 1.0$. Because the cut-edge $(u, v_{leaf})$ mathematically achieves this absolute global maximum, any naive global top-$k$ thresholding operation will prioritize this leaf-edge over any structural internal edge $(x,y)$ where $R(x,y) < 1.0$. Thus, $k$ leaf nodes will perfectly saturate a budget of $k$ edges, completely blinding the metric. $\blacksquare$

\textbf{Proof of Theorem 4 (Core-Shielded Resistance):} The $k$-core decomposition for $k \ge 2$ iteratively removes all nodes of degree $< 2$. By definition, all tree-like appendages (including the degree-1 leaves that cause $R=1.0$ saturation) are pruned. Since electrical current does not flow into dead-end branches when $s, t \in V_{core}$, the effective resistance $R(s,t)$ between any two nodes in the core is independent of the removed trees. Thus, the relative ordering of $R(x,y)$ for all core edges is perfectly preserved. The computational complexity of computing the pseudoinverse $L_{core}^+$ scales as $O(|V_{core}|^3)$, while the prerequisite $k$-core decomposition \citep{Batagelj2003} operates in linear $O(|V| + |E|)$ time, strictly bounding the total overhead. $\blacksquare$

\section{Experimental Details and Dataset Statistics}
\label{app:hyperparams}

\subsection{Dataset Statistics}
To rigorously evaluate the empirical claims and the homophily-based predictive diagnostic, we benchmarked CSER across a wide suite of standard datasets representing different structural regimes. Table~\ref{tab:datasets} details the exact properties of these graphs.

\begin{table}[h]
\centering
\begin{tabular}{lcccccc}
\hline
\textbf{Dataset} & \textbf{Type} & \textbf{Nodes} & \textbf{Edges} & \textbf{Features} & \textbf{Classes} & \textbf{Homophily ($h$)} \\
\hline
Cora & Citation & 2,485 & 5,069 & 1,433 & 7 & $\approx 0.81$ \\
CiteSeer & Citation & 2,120 & 3,679 & 3,703 & 6 & $\approx 0.74$ \\
Texas & WebKB & 183 & 295 & 1,703 & 5 & $\approx 0.11$ \\
Cornell & WebKB & 183 & 280 & 1,703 & 5 & $\approx 0.30$ \\
Wisconsin & WebKB & 251 & 466 & 1,703 & 5 & $\approx 0.21$ \\
\texttt{ogbn-arxiv} & Citation & 169,343 & 1,157,799 & 128 & 40 & $\approx 0.65$ \\
\hline
\end{tabular}
\caption{Summary statistics of the datasets utilized in the empirical benchmarking. Homophily refers to the edge homophily ratio.}
\label{tab:datasets}
\end{table}

\subsection{Hyperparameter Grids and Hardware Setup}
For all architectural ablations and end-to-end evaluations, we defined a strict, reproducible hyperparameter search space. 
For standard GCN and GraphSAGE models, we searched learning rates across $\{0.01, 0.005, 0.001\}$ and fixed the hidden dimension to 64. A dropout rate of $p=0.5$ is applied after the initial ReLU activation. 
For Graph Attention Networks (GAT), we required a distinct configuration to ensure stable attention head training: learning rate fixed to $0.005$, 8 attention heads in the first layer (hidden dimension 16 per head), and 1 attention head in the output layer, with a higher dropout of $p=0.6$ applied across the network.
All models are optimized using Adam with a weight decay of $5 \times 10^{-4}$. We train for exactly 200 epochs and record the test accuracy at the epoch that achieved the maximum validation accuracy.

The \texttt{ogbn-arxiv} wall-clock benchmarking was conducted on a standard local CPU architecture without reliance on distributed clusters or GPUs, empirically verifying the exact linear $O(|V| + |E|)$ complexity of the $k$-core topological filter under standard consumer hardware conditions.

\end{document}
